\documentclass[10pt,a4paper]{article}
\usepackage[margin=20mm]{geometry}
\usepackage{fontspec}
\usepackage{xetexko}
\setmainfont{Helvetica Neue}[BoldFont={Helvetica Neue Bold},ItalicFont={Helvetica Neue Italic},BoldItalicFont={Helvetica Neue Bold Italic}]
\setmainhangulfont{Apple SD Gothic Neo}[BoldFont={Apple SD Gothic Neo Bold}]
\usepackage{amsmath,amssymb}
\usepackage{graphicx}
\usepackage{xcolor}
\usepackage{caption}
\usepackage{enumitem}
\definecolor{revblue}{HTML}{1F4FD8}
\newcommand{\rev}[1]{{\color{revblue}#1}}
\renewcommand{\figurename}{그림}
\renewcommand{\tablename}{표}
\captionsetup{font=small,labelfont=bf,labelsep=period}
\linespread{1.25}
\setlength{\parskip}{3pt}

\begin{document}
\begin{center}
{\Large\bfseries 축소되는 서식지에서의 생존: 대규모 배틀로얄 텔레메트리 기반\\ 단계 조건부 탈락 위험 분석\par}
\vspace{4pt}
{\small 검토용 한글 초안 (서론–방법론) · 2026년 10월 7일 · \rev{파란 글씨}는 이번에 고친 부분\par}
\end{center}

\section{서론}
\label{sec:introduction}

% P1: 맥락과 관측 수단
사람, 차량, 동물처럼 움직이는 다수의 에이전트는 대규모 시공간 궤적 데이터를 남긴다~\cite{Zheng15,Atluri18}.
이러한 궤적을 사건 및 환경 기록과 결합하면 시간에 따른 에이전트의 위치뿐 아니라 에이전트 사이의 상호작용과 환경의 변화를 함께 분석할 수 있다.
특히 활동 영역이 시간에 따라 변하고 생존 경쟁이 발생하는 환경에서는 어떤 에이전트가 살아남고 어떤 에이전트가 탈락하는지를 이해하기 위해 영역의 변화와 경쟁을 함께 고려해야 한다.
배틀로얄 게임은 이러한 환경을 대규모로 관측할 수 있게 해 준다.
PUBG: BATTLEGROUNDS(이하 PUBG)에서는 최대 100명의 플레이어가 1인, 2인 또는 4인 팀으로 하나의 맵에 투입되어 마지막 한 팀이 남을 때까지 경쟁한다.
경기가 진행되는 동안 안전구역은 정해진 단계에 따라 줄어들고, 그 밖에 있는 플레이어는 지속적인 피해를 입는다.
따라서 안전구역은 구역 피해 없이 활동할 수 있는 영역이며, 그 축소는 팀의 이동과 교전 조건을 변화시킨다.
PUBG는 약 10초 간격의 플레이어 위치 기록과 교전, 아이템, 구역 이벤트를 포함하는 경기 텔레메트리를 공개한다~\cite{PUBGAPI}.
맵과 모드별로 정해진 규칙의 경기가 반복되므로, 이 데이터는 줄어드는 영역 아래의 생존을 반복적으로 관찰할 수 있는 관측 환경이 된다.

% P2: 추정 과제와 활용 맥락
이러한 경기에서 안전구역의 단계 전환은 팀의 중요한 의사결정 시점이 된다.
새 단계가 시작되어 다음 안전구역이 공개되면, 각 팀은 구역이 줄어들기 전에 언제 어떤 경로로 이동할지, 즉 로테이션을 정하고 현재 위치를 지킬지 교전할지를 판단해야 한다.
이러한 판단을 분석할 때 필요한 양은 각 단계 경계에서 생존한 팀이 이어지는 한 단계 안에 완전히 탈락할 조건부 확률, 즉 다음 단계 탈락 위험이다.
\rev{팀은 1인, 2인 또는 4인으로 구성되므로, 솔로 경기에서는 이 양이 곧 개인의 탈락 위험이다.}
이 양은 e스포츠에서도 활용될 수 있다.
프로 팀의 코치와 분석가에게는 단계별 로테이션 당시 팀이 처했던 위험을 복기하는 보조 지표가 될 수 있다.
여러 팀이 넓은 맵에 흩어져 동시에 움직이고 교전하는 대회 경기에서는 중계 옵저버가 관찰할 팀의 우선순위를 정하는 데 활용할 수 있다.
또한 개발사가 단계별 대기 시간, 축소 시간, 구역 밖 피해를 조정하려면 각 단계에서 관측되는 탈락 압박을 파악할 필요가 있다~\cite{PUBGBlueZone2026}.
본 연구는 공개 경기 기록을 활용한 관전자 관점의 위험 분석을 다루며, 특정 팀이 당시 알고 있던 정보만을 사용하는 의사결정 지원과는 구분된다.
반면 경기 단위의 우승확률은 경기 전체의 최종 결과를 요약하지만, 현재 생존한 팀이 이어지는 한 단계 안에 탈락할 확률을 직접 나타내지는 않는다.

% P3: 난점과 기존 설계의 한계
그러나 다음 단계 탈락 위험을 예측하기는 어렵다.
위험은 줄어드는 영역의 기하와 주변 팀과의 상호작용에 함께 달려 있고, 그 크기는 단계마다 달라진다.
또한 위치와 사건은 서로 다른 주기로 기록되므로 시간 정보를 정합적으로 맞춰야 한다.
기절한 플레이어와 탈락한 플레이어를 구분해야 하며, 단계 경계 이후의 정보가 입력에 포함되면 결과가 누설된다.
기존 PUBG 예측 연구는 생존과 경기 결과의 예측을 서로 다른 방식으로 다룬다.
Sagong 등~\cite{Sagong25a,Sagong25b}은 단계별 활동 점수를 이용하여 개별 플레이어의 이후 생존 여부를 예측한다.
KRAFTON의 예측기~\cite{KRAFTON}는 팀의 행동, 상태, 공간 정보를 팀 간 Transformer에 입력하고, Cox 계열의 학습 목적을 이용하여 미래 탈락 순서에 관한 상대위험 점수를 학습한 뒤 이를 우승확률로 변환한다.
그러나 위험집합 내 상대위험 점수와 정해진 다음 단계 안의 탈락확률은 서로 다른 추정 대상이다.
이러한 상대위험 학습 목적만으로는 팀 전체에 공통적인 단계별 위험 수준이 직접 식별되지 않으므로, 예측 구간과 단계별 기저위험을 명시적으로 다루는 설계와의 비교가 필요하다.
더욱이 두 기존 설계는 \rev{분석 단위(개인과 팀)}, 데이터, 평가 지표가 모두 달라서, 보고된 성능을 직접 비교하는 것만으로는 어떤 설계 선택이 예측 성능을 좌우하는지 알기 어렵다.

% P4: 생태학적 관점과 적용
줄어드는 영역 안에서의 생존은 서식지 변화 아래에서 집단의 존속을 다루는 생태학적 문제와 맞닿아 있으며, 서식지 손실은 생물다양성에 크고 일관된 부정적 영향을 미친다~\cite{Fahrig03}.
본 연구는 생태학의 모형을 그대로 옮기는 대신, 환경 조건에 따른 집단의 존속과 기간별 소멸확률이라는 관점을 PUBG에 적용한다(\ref{sec:related:survival}절).
안전구역의 축소를 구역 피해 없이 활동할 수 있는 영역의 단계적 감소로, 팀의 탈락을 해당 집단의 존속 종료로 대응시킨다.
이를 바탕으로 다음 단계 탈락 위험을 단계를 시간 단위로 하는 조건부 이산시간 hazard로 정식화한다.
이동을 환경과 연결하는 관점~\cite{Nathan08}도 함께 참고하여, 각 팀의 상태를 구역 경계에 대한 노출, 다음 구역에 도달하기 위한 이동 부담, 주변의 경쟁 밀도, 남은 팀원 수로 구성한 서식지 맥락으로 표현한다.
이 설계는 단계별 탈락이라는 추정 대상에 맞추어 위험의 시간 단위와 환경적 맥락을 함께 구성한다.

% P5: 연구 질문과 접근
이 관점은 세 가지 질문을 제기한다.
\begin{itemize}
    \item \textbf{RQ1:}
    서식지 축소 관점에서 구성한 구역 경계에 대한 노출, 이동 부담, 주변 경쟁 밀도, 남은 팀원 수는 다음 단계 탈락 위험과 어떤 방향과 크기의 관련을 보이며, 그 관계는 단계와 코호트에 걸쳐 얼마나 일관되는가?

    \item \textbf{RQ2:}
    단계 조건부 hazard 목적과 서식지 맥락 표현을 결합한 설계는 동일한 데이터와 평가 조건에서 기존 설계의 적용판과 어떤 예측 성능 차이를 보이는가?

    \item \textbf{RQ3:}
    입력 표현, 학습 목적, 단계별 기저위험, 모델 구조를 각각 변경할 때 예측 성능은 얼마나 달라지며, 그 차이는 코호트에 따라 어떻게 달라지는가?
\end{itemize}

이 질문에 답하기 위해 맵, 시점, 팀 규모로 나눈 12개 코호트의 공개 경기 13,948개를 사용한다.
이를 위해 약 10초 간격의 위치 기록과 비주기적 사건 로그를 결합하여, 예측 시점 이후의 정보가 입력되지 않는 팀--단계 학습 표본을 구성하는 파이프라인을 구축한다\rev{(그림~\ref{fig:pipeline})}.
이 파이프라인은 위치와 사건 기록을 공통 시간축에 정렬하고, \rev{생존 여부와 팀 간 관계를 포함한 모든 입력을 예측 시점 직전까지의 기록만으로 구성하며}, 탈락 라벨을 위치 스냅샷이 아닌 사건 기록으로 구성한다.
이어서 두 기존 설계를 동일한 팀--단계 표본에서 비교할 수 있도록 적용하고, \rev{각 적용판이 원래 과제에서도 제대로 작동하는지 함께 확인한다}.
다른 조건을 고정한 상태에서 입력 표현, 학습 목적, 단계별 기저위험, 모델 구조를 한 번에 하나씩 바꾸어 각 요소의 조건부 기여를 평가한다.

% 연구 질문에 대응하는 기여
본 논문의 기여는 세 질문에 대한 답으로 정리된다.
\begin{itemize}
    \item \textbf{RQ1에 대한 답.}
    다음 단계 탈락 위험을 단계적 서식지 축소 관점에서 정식화하고, 서식지 맥락과 단계 조건부 이산시간 hazard를 결합한 해석 가능한 모델을 제시한다.
    이 모델에서 이동 부담과 주변 경쟁 밀도는 탈락 위험을 높이고 남은 팀원 수는 위험을 낮추며, 세 공변량의 추정된 효과 방향은 평가 가능한 모든 코호트에서 일관된다.
    다만 효과의 크기는 코호트마다 다르다.
    구역 경계 안쪽으로의 깊이가 갖는 보호 효과는 안전구역이 줄어들수록 커져 중반 단계 이후 모든 코호트에서 관측되는 반면, 주변 경쟁 밀도의 효과는 단계가 진행될수록 작아진다.

    \item \textbf{RQ2에 대한 답.}
    이 정식화에 맞춘 설계, 즉 단계 조건부 hazard 목적과 서식지 맥락 표현을 두 기존 설계의 적용판과 대비하여 제시한다.
    동일한 표본과 평가 조건에서 예측 성능을 비교한 결과, 본 연구의 설계는 평가한 모든 코호트에서 두 적용판보다 우수한 점추정치를 보이며, \rev{특히 단계마다 달라지는 탈락 위험의 수준을 상대위험 기반 설계보다 정확하게 나타낸다.
    반대로 우승 팀 예측이라는 원래 과제에서는 상대위험 기반 설계가 더 정확하여, 두 학습 목적이 서로 다른 추정 대상에 맞추어져 있음을 보인다}.

    \item \textbf{RQ3에 대한 답.}
    통제된 비교를 통해 성능 차이에 대한 각 설계 요소의 조건부 기여를 평가한다.
    \rev{상대위험 기반 설계와의 차이는 단계별 기저위험의 결합과 다음 단계 탈락을 직접 학습하는 목적이 대부분 설명하며, 같은 학습기에서 본 연구의 입력과 재구성한 KRAFTON 입력의 예측력은 실질적으로 같고 두 입력을 결합하면 더 개선된다.}
    \rev{활동 점수 설계와의 차이는 입력 표현이 설명하며, 그 개선의 대부분은 서식지 맥락을 추가할 때 나타난다.}
    \rev{팀 축, 단계 축, 관계 편향과 같은 구조 요소의 추가 이득은 실질적 동등 범위 안에 머물며, 학습 예산을 늘려도 달라지지 않는다.}
\end{itemize}


\begin{figure}[t]
\centering
\includegraphics[width=\linewidth]{../../../paper/bigcomp2027/figures/fig_pipeline_clean.pdf}
\caption{데이터 파이프라인(위, \ref{sec:data}장)과 모델링·평가 절차(아래, \ref{sec:method}장).
예측 시점은 각 단계의 끝($b_k$, 다음 안전구역이 공개되는 시각)이며, 모든 입력은 그보다 10초 앞선 절단 시각($a_k=b_k-10$초) 이전의 기록으로 만든다.
라벨은 그 팀이 다음 단계 안에 탈락하면 1이다($Y=\mathbb{1}[b_k<T\le b_{k+1}]$).
각 설계는 자기 학습 방식을 유지하되 같은 경기 분할과 같은 시험 행에서 평가하며, 코호트별 차이는 확률효과 모형으로 통합한다.
실선은 데이터, 점선은 라벨의 흐름이다.}
\label{fig:pipeline}
\end{figure}

\section{관련 연구}
\label{sec:related}

\subsection{배틀로얄과 e스포츠의 결과 예측}
\label{sec:related:esports}

% 경기 종료 후의 순위 추정과 경기 중 예측의 구분
PUBG의 결과 예측 연구는 경기 종료 후의 요약 통계로 최종 순위를 추정하는 접근에서 출발하였으며~\cite{KagglePUBG}, 같은 유형의 통계로 순위를 예측할 때 유효한 요인과 모델이 경기 모드에 따라 달라짐이 보고되었다~\cite{LeeLee25}.
이러한 접근은 입력이 경기 종료 시점에 완성되므로, 경기 도중에 이용 가능한 정보로 이후의 위험을 예측하는 과제와는 구별된다.
다른 장르의 e스포츠에서는 경기 중에 갱신되는 정보를 이용한 예측이 연구되었다.
MOBA 장르에서는 실시간 승리 예측~\cite{Hodge21}과 시간 구간별 팀 간 상대적 우위의 평가~\cite{Yu19}가 제안되었고, 개인 단위의 짧은 구간 사망 예측도 Dota~2~\cite{Katona19,Ringer23}와 Counter-Strike~\cite{Marshall22}에서 다루어졌다.

% 배틀로얄의 가까운 선행연구와 본 연구의 비교 범위
배틀로얄의 경기 중 예측으로는 Sagong 등이 단계별 활동 및 교전 점수와 플레이 유형으로 개별 플레이어의 생존을 예측하였고~\cite{Sagong25a}, 단계별 전략 시퀀스로 승리를 예측하였다~\cite{Sagong25b}.
KRAFTON의 예측기는 팀의 행동, 상태, 공간 정보를 팀 간 Transformer에 입력하고, 생존분석을 이용하여 팀 단위의 실시간 우승확률을 추정한다~\cite{KRAFTON}.
공개 구현에는 이산시간 hazard(LogisticHazard)를 포함한 여러 생존모형 기준선과 팀 간 attention을 제거한 변형이 포함되어 있다~\cite{KRAFTON}.
이들 연구는 최종 순위와 우승확률, 개인의 생존 여부 등 서로 다른 분석 단위와 예측 구간을 다룬다.
본 연구는 안전구역의 단계 경계를 예측 시점으로 하고, \rev{솔로 경기의 개인을 1인 팀으로 포함하여} 현재 생존한 팀이 이어지는 한 단계 안에 완전히 탈락할 조건부 확률에 초점을 맞춘다.
이를 위해 활동 점수 기반 설계와 상대위험 기반 설계를 동일한 팀--단계 표본과 평가 조건에 맞추어 적용하고, 입력 표현, 학습 목적, 단계별 기저위험, 모델 구조가 예측 성능에 미치는 영향을 구분하여 평가한다.

\subsection{생존분석과 생태학적 위험 모델링}
\label{sec:related:survival}

% 상대위험, 절대적 사건확률 및 이산시간 모형
생존분석은 사건이 일어나기까지의 시간을 분석하면서, 사건이 관측되기 전에 추적이 종료되는 검열을 함께 다룬다.
Cox 비례위험모형은 기저위험의 함수 형태를 사전에 특정하지 않고 부분우도를 이용하여 공변량의 상대위험을 추정한다~\cite{Cox72}.
상대위험 점수 자체는 사건확률이 아니지만, 추정한 기저위험을 결합하면 생존확률과 구간별 조건부 사건확률을 산출할 수 있다.
이산시간 생존모형은 시간을 구간으로 나누고, 각 구간의 시작까지 생존한 대상이 해당 구간 안에 사건을 겪을 조건부 확률, 즉 구간별 hazard를 모형화한다~\cite{SingerWillett93}.
이때 구간별 기저위험을 통해 시간에 따른 공통적인 위험 수준의 변화를 반영할 수 있다.
신경망 기반 생존모형에서도 구간별 hazard를 학습하는 접근~\cite{Gensheimer19,Ren19}과 사건시간의 확률분포를 학습하는 접근~\cite{Lee18}이 제안되었으며, 시간에 따라 갱신되는 관측 이력을 이용한 동적 예측으로 확장되었다~\cite{Lee20}.

% 단계 경계에서의 동적 예측과 학습 목적 비교
예측 시점마다 그 시점의 생존자를 대상으로 당시까지 관측된 정보를 이용하여 이후의 위험을 추정하는 landmarking은 이러한 동적 예측을 구성하는 방법이다~\cite{vanHouwelingen07}.
본 연구는 각 단계 경계를 landmark로 삼아 관측 이력을 갱신하고, 이어지는 한 단계의 팀 탈락확률을 이산시간 hazard로 추정한다.
따라서 비교의 초점은 Cox 모형이 구간별 위험을 예측할 수 있는지 여부가 아니라, 상대위험 점수를 학습하는 목적과 다음 단계 탈락확률을 직접 학습하는 목적이 동일한 예측 과제에서 어떤 차이를 보이는지에 있다.
이를 위해 입력 표현과 공통 모델 구조를 고정한 상태에서 학습 목적을 비교하고, 단계별 기저위험의 기여를 별도로 평가한다.
Cox 계열 목적으로 학습한 모델에도 학습 집합에서 추정한 단계별 기저위험을 결합하여~\cite{PrenticeGloeckler78}, 기저위험의 유무가 아니라 학습 목적 자체의 차이를 비교한다.

% 생태학적 관점이 공변량 구성에 제공하는 근거와 적용 범위
생태학에서는 서식지의 양과 공간적 배치가 개체군의 존속과 어떻게 연결되는지를 다루어 왔다~\cite{Fahrig03,Hanski94}.
동적 점유 모형은 이전 기간에 점유된 장소에서 다음 기간에 종이 사라질 조건부 확률을 추정하며, 정착과 불완전한 탐지도 함께 모형화한다~\cite{MacKenzie03}.
이러한 접근은 환경 조건에 따른 기간별 존속 상태의 전이를 다룬다는 점에서 팀 탈락 문제와 연결되지만, 고정된 장소의 점유와 재정착을 다루는 모형을 이동하는 팀에 그대로 적용할 수는 없다.
또한 이동생태학은 이동을 개체의 상태와 능력, 외부 환경이 함께 작용하는 과정으로 해석한다~\cite{Nathan08}.
본 연구는 이들 관점을 참고하여 안전구역의 축소 아래에서 팀의 존속과 관련될 수 있는 구역 경계에 대한 노출, 이동 부담, 주변 경쟁 밀도, 남은 팀원 수를 서식지 맥락으로 구성한다.
생태학적 관점은 위험 공변량을 구성하고 그 관계를 해석하는 근거로 사용하며, 공변량과 탈락 위험의 관련성 및 그 일관성은 경기 자료에서 검증한다.

\subsection{시공간 다중 에이전트 모델링}
\label{sec:related:multiagent}

% 상호작용과 시간적 이력을 표현하는 구조
시공간 데이터 마이닝은 이동 궤적과 사건을 이용한 예측, 군집, 이상 탐지 등을 다루어 왔으며, 딥러닝을 이용하여 시공간 의존성을 학습하는 방법도 연구되었다~\cite{Atluri18,Wang22}.
다수의 에이전트가 상호작용하는 장면에서는 에이전트 축과 시간 축의 attention을 분리하여 결합하고, 이를 통해 에이전트 사이의 관계와 시간에 따른 변화를 함께 표현하는 방법이 제안되었다~\cite{Ngiam22,Capellera24}.
그래프 구조를 attention에 반영하는 접근으로는 노드 간 최단경로 거리와 간선 정보를 attention의 편향으로 인코딩하는 Graphormer가 있다~\cite{Ying21}.
스포츠와 e스포츠에서도 선수 간 관계를 그래프로 표현하여 미식축구의 전진 거리와 Counter-Strike의 라운드 승리와 같은 경기 내 결과를 예측한 연구가 있다~\cite{Xenopoulos21}.

% 서식지 표현과 hazard 목적 위에서의 구조적 추가 이득
이러한 구조는 에이전트 간 상호작용과 시간적 이력을 학습할 수 있는 표현력을 제공하지만, 입력 정보, 학습 목적, 모델 구조가 동시에 달라지는 비교에서는 관측된 성능 차이를 구조 자체의 기여로 해석하기 어렵다.
본 연구는 팀 축 attention, 단계 축 attention, 관계 편향을 갖춘 모델을 참조 모델로 사용하되, 서식지 맥락 표현과 단계 조건부 hazard 목적을 고정한 상태에서 각 구조 요소를 한 번에 하나씩 추가하여 그 조건부 기여를 평가한다.

\section{데이터}
\label{sec:data}

\subsection{PUBG 경기 텔레메트리}
\label{sec:data:telemetry}

% P1: 무엇이 기록되는가
PUBG는 공개 API를 통해 경기마다 텔레메트리를 제공한다~\cite{PUBGAPI}.
텔레메트리는 경기에 참가한 모든 플레이어의 위치를 약 10초 간격으로 기록하며, 표본으로 추출한 약 1,738만 개의 연속 위치 간격에서 중앙값은 10초, 11초 이하의 비율은 99.93\%였다.
위치 기록과 함께 피해, 기절, 처치, 아이템의 획득·장착·사용, 안전구역의 상태 변화, 단계 전환, 착지와 같은 사건이 발생 시각과 함께 기록된다.
따라서 이 자료는 모든 팀의 이동 궤적과 팀 사이의 교전, 그리고 안전구역의 변화를 하나의 경기 안에서 함께 담는다.

% P2: 수집 규모와 희소성
본 연구는 PUBG 공개 API\rev{의 표본 경기 목록을 하루 단위로 조회하여} 수집한 스팀 서버의 경기 가운데 \ref{sec:data:benchmark}절의 기준으로 선정한 13,948경기를 분석에 사용한다.
이 경기들은 2026년 3월 26일부터 6월 2일까지 생성되었으며, 위치 표본 약 9,643만 건과 사건 약 2억 2,597만 건, 모두 약 3억 2,240만 건의 기록을 포함한다(표~\ref{tab:corpus}).
경기 종료 후의 통계를 집계한 공개 데이터셋은 존재하지만~\cite{KagglePUBG}, 경기 중 모든 플레이어의 위치를 이 해상도로 담은 자료는 드물며, 기존 배틀로얄 예측 연구가 사용한 경기 내 기록도 공개되지 않았다~\cite{Sagong25a,KRAFTON}.

\begin{table}[t]
\centering
\caption{분석에 사용한 13,948경기의 텔레메트리 기록 유형별 건수}
\label{tab:corpus}
\begin{tabular}{lr}
\hline
기록 유형 & 건수 \\
\hline
위치 표본 (약 10초 간격) & 96,426,398 \\
피해 & 90,178,937 \\
아이템 획득·버림 & 67,865,634 \\
아이템 장착 & 46,218,920 \\
아이템 사용 & 15,229,394 \\
안전구역 상태 갱신 & 2,134,448 \\
착지 & 1,673,902 \\
처치 & 1,267,179 \\
기절 & 1,199,290 \\
단계 전환 & 200,979 \\
\hline
합계 & 322,395,081 \\
\hline
\end{tabular}
% DB 확인: 벤치마크 match_id 13,948개로 COUNT (2026-10-06).
\end{table}

\subsection{학습 표본 구성}
\label{sec:data:samples}

% P3: 공통 시간축, 단계 경계, 인과적 절단
원시 기록을 학습 표본으로 바꾸려면 세 가지 시간 처리가 필요하다\rev{(그림~\ref{fig:pipeline}의 B)}.
첫째, 모든 기록을 하나의 시간축에 놓아야 한다.
위치 기록에 함께 저장된 경기 경과 시간은 사건 기록의 타임스탬프와 기준시각이 달라 경기별로 최대 41.9초까지 어긋나므로, 본 연구는 모든 기록의 타임스탬프를 경기 시작 사건으로부터의 경과 시간으로 변환하여 하나의 시간축으로 사용한다.
둘째, 단계 경계는 단계 전환 사건으로 정한다.
텔레메트리는 새 안전구역의 공개와 구역 축소의 시작을 단계 전환 사건으로 기록하며, 한 구역으로의 축소가 끝나는 시각에 다음 구역이 공개된다~\cite{PUBGAPI}.
본 연구는 $k$번째 안전구역이 공개된 시점부터 그 구역으로의 축소가 끝날 때까지를 단계 $k$로 정의하고, 축소가 끝나 다음 안전구역이 공개되는 시각을 단계 $k$의 종료 시각 $b_k$로 둔다.
셋째, 예측 시점은 $b_k$로 두고, 모든 입력은 그보다 10초 앞선 특징 절단 시각 $a_k=b_k-10$초 이전의 기록만으로 구성한다.
이 제한은 특징뿐 아니라 생존 인원, 팀 사이의 관계, attention 마스크에도 적용되며, 생존 여부 역시 최종 탈락 시각이 아니라 $a_k$ 이전의 처치와 위치 기록으로 판정한다.
처치된 플레이어는 처치 후 20초 이상 지난 시점에 체력이 남은 위치 기록이 나타날 때에만 재투입된 것으로 본다.
12개 코호트에서 17경기씩 고른 204경기의 1,293개 단계 경계에서 $a_k$ 이후의 원시 기록을 모두 지우고 입력을 다시 구성한 결과, 모든 입력이 원래 입력과 일치하였다.
특징 절단 시각에 생존한 플레이어의 99.4\%는 직전 15초 안에 위치 표본을 가진다.

% P4: 사건 기반 탈락 라벨
탈락 라벨은 사건 기록으로 정의한다.
기절은 이후 소생이 가능한 상태이므로 최종 처치와 구분하며, 팀의 탈락 시각은 재투입되지 않은 마지막 팀원이 처치된 시각으로 정의한다.
이렇게 정한 탈락 시각은 모든 탈락 팀에서 게임이 팀의 최종 순위를 기록한 처치 사건의 시각과 일치하였다.
우승 팀은 탈락 사건 없이 경기가 끝나므로 경기 종료 시각에서 관측을 종료한다.
단계 $k$의 위험집합은 예측 시점 $b_k$에 생존한 팀들이며, 각 팀의 라벨은 다음 단계의 종료 시각 $b_{k+1}$까지 탈락하면 1, 그렇지 않으면 0이다.
다음 단계가 경기의 마지막 단계이면 $b_{k+1}$ 대신 경기 종료 시각을 사용한다.
경기 종료 기록이 없거나 라벨 구간에서 위치 기록이 60초 넘게 끊긴 경기는 없었다.
예측 시점은 단계 1부터 7까지의 경계에 둔다.
따라서 분석의 단위는 경기, 팀, 단계 경계의 조합인 팀--단계 표본이다.
\rev{솔로 경기에서는 팀이 한 명의 플레이어로 이루어지므로, 팀--단계 표본이 곧 플레이어의 단계별 표본이다.}

\subsection{벤치마크}
\label{sec:data:benchmark}

% P5: 모집단 정제와 코호트
분석 대상은 일반 경기와 랭크 경기로 한정하였다.
실제 경쟁이 이루어지지 않는 튜토리얼과 훈련장 경기는 제외하였고, 일반적인 경기 양상과 다른 두 모집단도 제외하였다.
계정의 약 92\%가 AI인 airoyale 경기와, 경기마다 약 30개의 비플레이어 전투 계정이 기록되는 Vikendi 맵의 경기이다.
일반 경기에도 AI 계정이 섞여 있으므로, 각 팀의 AI 계정 비율을 입력 특징에 포함하였다.
경기는 맵, 시점(3인칭·1인칭), 팀 규모(1인·2인·4인)의 조합으로 코호트를 나누었으며, airoyale 경기를 제외하기 전 기준으로 300경기 이상인 조합을 코호트로 정하였다.
최종 벤치마크는 12개 코호트의 13,948경기이며, 이 중 11,285경기는 일반 경기, 2,663경기는 랭크 경기이다.
\rev{팀 규모별로는 솔로 1개, 듀오 2개, 스쿼드 9개 코호트이다.}
% 보류: 랭크 경기와 대회의 안전구역 설정 일치는 규정 원문 또는 경기 설정 기록으로 확인한 뒤 복원.
% 랭크 경기는 e스포츠 대회와 같은 안전구역 설정을 사용한다~\cite{PUBGRules}.

% P6: 분할과 표본 수
각 코호트의 경기는 경기 단위로 학습, 검증, 시험 집합에 8:1:1로 나누었으며, 분할은 고정하여 모든 비교에 같은 분할을 사용한다.
그 결과 학습, 검증, 시험 집합은 각각 11,188경기, 1,376경기, 1,384경기이며, 코호트의 크기는 224경기부터 4,159경기까지이다.
위험집합에 속한 팀--단계 표본은 모두 817,766개이고, 그중 217,804개(26.6\%)가 다음 단계 안에 탈락하였다.
\rev{표본의 83.1\%는 스쿼드, 13.0\%는 듀오, 3.9\%는 솔로 코호트에서 나온다.}
다음 단계 탈락 비율은 단계 1의 19.6\%에서 단계 7의 50.4\%로 높아진다.

\section{방법론}
\label{sec:method}

\rev{그림~\ref{fig:pipeline}의 C와 D는 이 장의 설계와 평가 절차를 요약한다.}

\subsection{문제 정식화}
\label{sec:method:formulation}

% P1: 예측 시점, 특징 절단 시점, 위험집합, 라벨
경기 $m$에서 단계 $k$의 종료 시각을 $b_{m,k}$, 특징 절단 시각을 $a_{m,k}=b_{m,k}-\Delta$($\Delta=10$초)라 하자.
팀 $i$의 탈락 시각을 $T_{m,i}$, 경기 종료 시각을 $C_m$이라 하면, 관측되는 시각은 $\widetilde T_{m,i}=\min(T_{m,i},C_m)$이고 탈락 여부는 $\delta_{m,i}=\mathbb{1}[T_{m,i}\le C_m]$이다. 우승 팀은 $\delta_{m,i}=0$이다.
예측 시점은 단계 경계 $b_{m,k}$이며, 위험집합은 이 시점에 생존한 팀들의 집합 $\mathcal{R}_{m,k}=\{i:\widetilde T_{m,i}>b_{m,k}\}$이다.
위험집합의 각 팀에 대해 라벨은 $Y_{m,i,k}=\mathbb{1}[b_{m,k}<T_{m,i}\le b_{m,k+1}]$이며, 예측 대상은 다음 단계 탈락 위험
\begin{equation}
\lambda_{m,i,k}=\Pr\!\left(T_{m,i}\le b_{m,k+1} \mid T_{m,i}>b_{m,k},\,\mathcal{H}_{m}(a_{m,k})\right)
\label{eq:hazard}
\end{equation}
이다. 여기서 $\mathcal{H}_{m}(t)$는 시각 $t$까지의 경기 기록이다.
다음 단계가 경기의 마지막 단계이면 $b_{m,k+1}=C_m$으로 둔다.
위험집합의 구성에는 예측 시점의 생존 여부를 사용하지만, 모델의 입력 특징, 생존 인원, 팀 간 관계, attention 마스크는 특징 절단 시각 이전의 정보로만 구성한다.
이 정의는 각 단계 경계를 landmark로, 다음 단계를 예측 구간으로 하는 이산시간 hazard에 해당한다.
본 연구는 hazard를 로짓 척도에서 단계별 기저위험과 팀 상태의 함수로 모형화한다.
\begin{equation}
\mathrm{logit}\,\lambda_{m,i,k}=\alpha_k+f(\mathbf{x}_{m,i,k},\mathcal{G}_{m,k})
\label{eq:logit}
\end{equation}
여기서 $\mathbf{x}_{m,i,k}$는 팀의 서식지 맥락, $\mathcal{G}_{m,k}$는 같은 단계의 다른 팀들과의 관계이다.

% P2: 상대위험 학습 목적과의 차이
본 연구에서 비교하는 상대위험 학습 목적은 같은 위험집합 안에서 팀들의 점수 차이를 학습한다.
표준 Cox 부분우도의 각 항은 사건 팀의 점수가 위험집합에서 차지하는 비율 $\exp(s_i)/\sum_{j}\exp(s_j)$이며, KRAFTON의 weighted Cox는 이를 가중한 변형이다.
이러한 목적은 같은 시점의 모든 점수에 공통으로 더해지는 이동에 불변이므로, 그 자체로는 식~\eqref{eq:logit}의 $\alpha_k$와 같은 단계별 공통 위험 수준을 식별하지 못한다.
기저위험을 따로 추정하여 결합하면 구간별 확률을 산출할 수 있으므로, 본 연구는 상대위험 점수에 단계별 기저위험을 결합한 변형도 함께 비교한다(\ref{sec:method:decomposition}절).

\subsection{서식지 맥락 표현}
\label{sec:method:representation}

% P3: 구성개념 네 가지
서식지 축소 관점에서 각 팀의 상태를 네 가지 구성개념으로 정의한다.
첫째, 구역 경계에 대한 노출은 팀 중심과 안전구역 경계 사이의 부호 있는 거리로 측정하며, 구역 밖에 있는 깊이와 구역 안쪽으로 들어온 깊이로 나누어 사용한다.
둘째, 이동 부담은 안전구역에 도달하는 데 필요한 속도와 팀이 실제로 이동한 속도의 차이, 즉 이동 부족분으로 측정한다~\cite{Nathan08}.
필요 속도는 특징 절단 시각 직전의 팀 중심에서 그 시점에 공개된 안전구역 경계까지의 거리(구역 안이면 0)를 단계 길이로 나눈 값이고, 이동 속도는 단계 시작부터 특징 절단 시각까지의 이동 경로 길이를 같은 단계 길이로 나눈 값이다.
단계 길이에는 경기별 관측값 대신 학습 집합에서 코호트, 경기 유형, 단계별로 구한 중앙값을 사용하여, 특징 절단 시각 이후의 정보가 들어가지 않게 하였다.
셋째, 주변 경쟁 밀도는 팀 주변 200\,m 안에 있는 다른 팀의 수이다.
넷째, 남은 팀원 수는 특징 절단 시각에 생존한 것으로 관측된 팀원의 수이다.
학습 집합에서 필요 속도는 구역 밖 깊이와 순위 상관이 0.82로 높아 정보가 중복될 가능성이 있으므로, 실제 이동을 함께 반영하는 이동 부족분(상관 0.38)을 사용하였다.
이 선택은 결과를 확인한 뒤 이루어진 탐색적 특징 설계이다.

% P4: 전체 특징 표현과 관계
예측 모델의 입력은 이 구성개념을 포함하는 150차원의 팀--단계 특징이다.
특징은 경기 맥락, 팀 상태, 안전구역, 이동, 노출, 자원, 공간, 관계, 서식지의 아홉 계열로 구성되며, 여기에 팀 중심 좌표와 팀의 AI 계정 비율을 더한다.
같은 단계의 팀들 사이에는 세 유형의 관계를 정의한다.
근접 관계는 단계 시작부터 특징 절단 시각까지 두 팀의 평균 거리 $d$에 대해 가중치 $\exp(-d/200)$를 갖는 대칭 관계이다.
교전 관계는 같은 구간에서 한 팀이 다른 팀에 가한 피해, 기절, 처치를 결합한 강도 $\log(1+\text{강도})$를 가중치로 갖는 방향 관계이며, 교전이 있었던 팀 쌍에서는 근접 관계를 교전 관계로 대체한다.
공동 압박 관계는 두 팀이 모두 구역 밖에 있거나 경계 안쪽 150\,m 이내에 있으면서 서로 200\,m 이내에 있을 때의 대칭 관계이다.

\subsection{RQ1: 생태 hazard 모델}
\label{sec:method:rq1}

% P5: 해석 가능한 모델과 일관성 평가
RQ1에 답하기 위해 \ref{sec:method:representation}절의 구성개념만으로 식~\eqref{eq:logit}을 로지스틱 회귀로 추정한다.
\begin{equation}
\mathrm{logit}\,\lambda_{m,i,k}=\alpha_k+\sum_{c}\beta_c z_{c,m,i,k}
\label{eq:ehm}
\end{equation}
여기서 $z_c$는 학습 집합에서 표준화한 구성개념(구역 밖 깊이, 구역 안쪽 깊이, 이동 부족분, 주변 경쟁 밀도, 남은 팀원 수)이며, 전체 생존 팀 수의 로그를 통제 변수로 포함한다.
솔로 코호트에서는 남은 팀원 수가 상수이므로 제외한다.
모델은 코호트마다 학습 집합에서 추정하며, 표준오차는 경기 단위로 군집화하여 계산한다.
코호트별 계수는 DerSimonian--Laird 확률효과 모형으로 통합하고, 통합 추정치와 95\% 신뢰구간, 코호트 간 이질성($\tau^2$, $I^2$), 그리고 계수의 부호가 통합 추정치와 같은 코호트의 수를 보고한다.
단계에 따른 변화를 보기 위해 각 구성개념과 단계 구간(초반 1--2, 중반 3--5, 후반 6--7)의 상호작용을 추가한 모델을 추정하고, 구간별 기울기와 구간 간 기울기가 같은지에 대한 검정 결과를 보고한다.

\subsection{RQ2: 비교 설계}
\label{sec:method:rq2}

% P6: 본 연구의 설계
본 연구의 설계는 식~\eqref{eq:logit}의 단계 조건부 hazard를 다음 단계 탈락에 대한 교차엔트로피로 학습하고, \ref{sec:method:representation}절의 서식지 맥락 표현을 입력으로 사용한다.
설계가 특정 학습기에 의존하지 않음을 보이기 위해 두 학습기로 구현한다.
첫째는 단계 경계의 팀 특징에 대한 gradient boosting 분류기(HGB; 트리 200개, 학습률 0.06)로, 단계 번호를 입력 특징으로 포함하며 식~\eqref{eq:logit}의 $\alpha_k$를 별도의 항으로 두지는 않는다.
둘째는 \rev{구조 요소의 기여를 분해하기 위한} 참조 모델인 PTR-Former로, 팀--단계 토큰에 단계 임베딩을 더한 뒤 같은 팀의 단계 이력에 대한 인과적 attention(단계 축)과 같은 단계의 생존 팀들에 대한 attention(팀 축)을 번갈아 적용하며, 팀 축 attention의 로짓에는 관계 유형별로 학습되는 편향을 더한다.
PTR-Former는 다음 단계 hazard의 교차엔트로피에 최종 순위에 대한 쌍별 순위 손실을 0.1의 가중치로 더하여 학습한다\footnote{차원 128, attention 헤드 4개, 블록 2개, dropout 0.1을 사용하였다. AdamW(학습률 $10^{-3}$, 가중치 감쇠 $10^{-4}$)와 코사인 학습률 감소로 최대 12 epoch 동안 경기 8개씩 묶어 학습하였으며, 검증 집합의 로그 손실이 3 epoch 동안 개선되지 않으면 학습을 멈추고 가장 좋았던 시점의 가중치를 사용하였다.}.

% P7: 기존 설계의 적용
두 기존 설계는 원 설계의 입력과 학습 목적을 유지하면서 본 연구의 표본에 맞게 적용하였으며, 적용 내용은 표~\ref{tab:adapt}에 정리하였다.
Sagong 등의 설계~\cite{Sagong25a}는 플레이어 단위의 활동 지표로 구성되므로 2인·4인 코호트에서는 팀 단위로 집계하였고\rev{(솔로 코호트에서는 원 설계와 같은 개인 단위이다)}, 본 연구의 정규화 자료에 포함되지 않은 공격 시도와 어시스트 기록은 가장 가까운 사건으로 대체하였다.
KRAFTON의 예측기~\cite{KRAFTON}는 공개 구현(commit 99ed9d1)을 기반으로, 입력 지표를 본 연구의 텔레메트리로 재구성하고 공개 경기와 평가 과제에 맞춘 적용판으로 구성하였다.
모든 설계는 같은 경기 분할을 사용하며, 평가는 같은 팀--단계 위험집합과 다음 단계 탈락 라벨에서 수행한다.
RQ2에서는 각 적용판의 원래 학습 표집과 감독 목적을 유지하고, RQ3에서는 학습 표본과 입력을 추가로 통제한다.

% P7b: 적용판의 충실도 점검
적용판이 본 연구의 과제에 불리하게 구성되지 않았는지 확인하기 위해, 각 적용판을 원래 과제의 지표로도 평가한다.
KRAFTON 적용판은 원 구현이 정의한 우승 팀 정확도와 탈락 순서의 C-index로, Sagong 적용판은 \rev{분석 단위가 원 설계와 같은 솔로 코호트에서} 원 논문과 같은 개인 단위의 생존 분류로 평가한다.
본 연구의 설계도 같은 지표로 평가하여, 과제가 바뀔 때 설계 간 우열이 어떻게 달라지는지 확인한다.

\begin{table}[t]
\centering
\caption{기존 설계의 적용}
\label{tab:adapt}
\small
\begin{tabular}{p{1.3cm}p{3.1cm}p{3.1cm}}
\hline
항목 & Sagong 적용판 & KRAFTON 적용판 \\
\hline
원 설계 & 플레이어 단위, 다음 반 단계 생존 분류 & 팀 단위, 사망 시각 Cox 학습 후 우승확률로 변환 \\
입력 & 단계별 활동 지표 8개, FightScore, ActiveScore, 플레이 유형, 유형별 전이 생존율 (12개) & 팀 지표 55개(원 정의와 같음 32, 근사 22, 재구성하지 못해 상수 0으로 대체 1), 팀원 위치, 구역·맵 임베딩 \\
학습기 & XGBoost, Random Forest, 로지스틱 회귀 & 8층 Transformer, weighted Cox \\
학습·평가 시점 & 단계 경계 / 단계 경계 & 원 설계의 시점 표집 / 단계 경계 \\
적용 결정 & 공격은 적에게 가한 피해 사건 수로, 어시스트는 다른 플레이어의 처치 직전 60초 안의 피해로 대체; 2·4인 팀은 횟수 합, 밀집도 최댓값, 이동 거리 평균으로 집계; 반 단계 대신 단계를 집계 구간으로 사용 & 팀 수 상한 제거, 탈락 팀의 attention 차단, 평가 코드의 점수 방향 해석 수정(99ed9d1), 검증 집합에서 Platt 보정(같은 검증 집합에서 단계별 기저위험을 추정한 +H0 변형 병행) \\
설정 & 라이브러리 기본값(원 논문 미기재) & 원 설계의 설정 \\
\hline
\end{tabular}
\end{table}

\subsection{RQ3: 단일 요인 분해}
\label{sec:method:decomposition}

% P8: 분해 사다리
RQ3에 답하기 위해 다른 조건을 고정한 채 한 번에 하나의 설계 요인만 바꾸는 비교를 구성한다(표~\ref{tab:ladder}).
입력 표현의 효과는 학습기를 HGB로 고정하고 Sagong 등의 입력 또는 재구성한 KRAFTON 입력을 본 연구의 입력으로 바꾸어 측정한다.
Sagong 등의 입력과의 비교는 서식지 맥락을 먼저 추가한 뒤 나머지 입력을 바꾸는 두 단계로 나누어, 차이 가운데 서식지 맥락이 차지하는 부분을 구분한다.
학습 목적과 단계별 기저위험의 효과는 KRAFTON의 Transformer 구조에 본 연구의 입력을 넣고 weighted Cox 목적으로 학습한 모델(KR-Cox)에서 출발한다.
KRAFTON 원본에서 KR-Cox로 넘어가는 비교는 입력과 학습 표본이 함께 바뀌는 묶음 비교이므로, 분해 단계와 구분하여 보고한다.
KR-Cox의 학습 목적을 다음 단계 탈락의 교차엔트로피로 바꾼 모델을 KR-Hazard, KR-Hazard에 단계별 로짓 절편 $\alpha_k$를 추가한 모델을 KR-Hazard+$\alpha$라 한다.
KR-Cox와 KR-Hazard는 구조, 입력, 학습 경기와 절단 시점, 초기 가중치, 갱신 횟수가 같고 학습 목적만 다르다.
다만 KR-Cox는 특징 절단 시각에 생존한 팀들의 이후 탈락 순서를, KR-Hazard는 예측 시점의 위험집합에서 다음 단계의 탈락 여부를 감독 정보로 사용하므로, 이 비교는 원래의 상대위험 학습 목적과 과제에 맞춘 다음 단계 위험 학습 목적의 차이를 나타낸다.
입력에는 단계 번호가 포함되어 있으므로, KR-Hazard+$\alpha$와의 비교는 단계 정보를 포함한 입력을 유지한 상태에서 명시적인 단계별 로짓 절편을 추가하는 효과를 나타낸다.
Cox 점수 자체에는 단계별 위험 수준이 담기지 않으므로, KR-Cox의 점수에 학습 집합에서 추정한 단계별 기저위험을 결합한 KR-Cox+H0도 함께 평가한다~\cite{PrenticeGloeckler78}.
따라서 학습 목적의 비교는 KR-Cox에 기저위험을 결합하는 단계와, 기저위험을 갖춘 KR-Cox+H0를 KR-Hazard로 바꾸는 단계로 나누어 측정한다.
뒤의 단계는 기저위험을 갖춘 Cox 모델과 다음 단계 탈락확률을 직접 학습한 모델의 차이를 보여 준다.
아키텍처와 구조의 비교에서는 학습 목적을 맞추기 위해 순위 보조 손실을 사용하지 않는다.
아키텍처의 효과는 같은 학습 목적과 입력에서 KRAFTON 구조를 PTR-Former의 팀 축만 남긴 모델(PTR-Team)로 바꾸어 측정하고, 구조의 효과는 PTR-Team에 단계 축과 관계 편향을 추가한 PTR-Former와 비교하여 측정한다.
순위 보조 손실의 효과는 같은 구조에서 이를 추가한 모델과 비교하여 별도로 보고한다.
구조 요인은 단계 축과 팀 축을 2$\times$2로 켜고 끄는 요인 실험으로 평가하고, 관계의 대응을 경기 안에서 무작위로 섞어 다시 학습한 대조 모델로 관계 편향의 기여를 확인한다.
또한 안전구역의 원시 기하(현재·다음 구역의 중심과 반경)만을 갖는 구역 객체를 토큰으로 추가하여, 수작업 서식지 특징을 뺀 입력에서 그 정보를 대신 학습하는지 확인한다.
마지막으로 HGB의 입력에 재구성한 KRAFTON 지표를 추가하여 두 표현의 보완성을 확인한다.

\begin{table}[t]
\centering
\caption{단일 요인 분해의 단계}
\label{tab:ladder}
\small
\begin{tabular}{p{0.5cm}p{2.9cm}p{1.9cm}p{1.9cm}}
\hline
단계 & 비교 & 바꾸는 요인 & 고정하는 조건 \\
\hline
(a1) & HGB: Sagong 입력 $\rightarrow$ +서식지 맥락 & 서식지 맥락 & 학습기, 학습 목적 \\
(a2) & HGB: Sagong 입력+서식지 맥락 $\rightarrow$ 본 연구 입력 & 나머지 입력 & 학습기, 학습 목적 \\
(b) & HGB: KRAFTON 입력 $\rightarrow$ 본 연구 입력 & 입력 표현 & 학습기, 학습 목적 \\
(c1) & KR-Cox $\rightarrow$ KR-Cox+H0 & 단계별 기저위험 결합 & 구조, 입력, 표본, 학습 목적 \\
(c2) & KR-Cox+H0 $\rightarrow$ KR-Hazard & 학습 목적 & 구조, 입력, 표본 \\
(d) & KR-Hazard $\rightarrow$ KR-Hazard+$\alpha$ & 명시적 단계 절편 & 구조, 입력, 학습 목적 \\
(e) & KR-Hazard+$\alpha$ $\rightarrow$ PTR-Team & 아키텍처 & 학습 목적(순위 손실 없음), 입력 \\
(f) & PTR-Team $\rightarrow$ PTR-Former & 단계 축, 관계 편향 & 학습 목적(순위 손실 없음), 입력, 팀 축 \\
(g) & HGB: 본 연구 입력 $\rightarrow$ +KRAFTON 입력 & 입력 확장 & 학습기, 학습 목적 \\
\hline
\end{tabular}
\end{table}

\subsection{평가 프로토콜}
\label{sec:method:evaluation}

% P9: 지표
주 지표는 시험 집합의 위험집합 전체에 대한 AUC이다.
보완 지표로 두 가지를 함께 보고한다.
단계 안 AUC는 같은 단계 번호의 표본을 경기에 걸쳐 모아 단계별로 계산한 AUC의 평균으로, 단계마다 공통으로 더해지는 위험 수준의 차이를 제외한 판별력을 나타내며, 시험 표본이 30개 이상이고 두 라벨이 모두 있는 단계만으로 계산한다.
경기--단계 내부 AUC는 같은 경기의 같은 단계 위험집합 안에서 팀을 구분하는 판별력으로, 위험집합별 AUC를 양성--음성 쌍의 수로 가중하여 통합하고 라벨이 하나뿐인 위험집합은 제외한다.
로그 손실과 Brier 점수는 확률 예측의 전반적인 정확도를 나타낸다.
단계별 보정 오차는 단계마다 평균 예측확률과 관측 탈락률의 차이의 절댓값을 구해 표본 수로 가중 평균한 값으로, 설계가 단계마다 달라지는 위험 수준을 표현하는지를 보여 준다.
상대 순위에 대해서는 최종 순위와의 Spearman 상관과 상위 5위 판별 AUC를 사용한다.
상대 순위 지표에는 각 모델의 위험 점수를 사용한다. 순위 보조 손실로 학습한 PTR-Former는 순위 출력을, 나머지 모델은 다음 단계 탈락확률 또는 Cox 점수를 사용한다.

% P10: 통계적 추론과 분석의 구분
두 설계의 차이는 코호트마다 경기 단위로 2,000회 재표집한 쌍대 bootstrap으로 신뢰구간을 구하며, 이 구간은 학습된 모델을 고정한 상태에서 시험 표본이 달라질 때의 변동을 나타낸다.
12개 코호트의 결과는 DerSimonian--Laird 확률효과 모형으로 통합하고, REML과 Hartung--Knapp 방식의 결과를 민감도 분석으로 함께 확인한다.
의미 있는 최소 효과 크기는 AUC 0.005로 둔다.
통합 추정치의 95\% 신뢰구간 하한이 0보다 크고 점추정치가 0.005 이상이면 개선, 95\% 신뢰구간 상한이 0보다 작고 점추정치가 $-0.005$ 이하이면 저하, 90\% 신뢰구간이 $\pm0.005$ 안에 있으면 실질적 동등, 그 밖에는 불확정으로 분류한다.
개선 분류는 양의 차이에 대한 통계적 근거와 점추정치의 크기를 함께 요구하는 기준이며, 실제 효과가 0.005 이상임을 입증하는 검정은 아니다.
구조 요인의 다섯 비교는 실험 전에 타임스탬프가 남는 커밋으로 분석 계획을 고정한 사전 지정 분석이며, 다섯 개의 통합 검정에 Holm 방법을 적용한다\footnote{분석 계획은 2026년 10월 5일에 고정하였다. 다섯 비교는 단계 축과 팀 축의 2$\times$2 요인 실험에서 다른 축의 두 수준에 대해 평균한 단계 축의 주효과(AUC), 같은 방식의 팀 축 주효과(AUC), 팀 축 주효과의 Spearman 상관 차이, 관계의 대응을 경기 안에서 섞은 대조 모델 대비 관계 편향의 효과(AUC), 서식지 특징을 뺀 입력에서 구역 객체 토큰의 효과(AUC)이다. 계획을 고정한 뒤 데이터 처리의 누수를 발견하여 고쳤으며, 모든 결과는 고친 자료에서 다시 얻었다. 이 수정은 이전 결과를 확인한 뒤 결정하였다.}.
% 익명 심사본에서는 커밋 해시를 생략하고, 최종본에 추가한다(f59bac9).
생태 hazard 모델, 기존 설계와의 비교, 적용판의 원래 과제 평가, 단일 요인 분해, 단계 안 AUC, 경기--단계 내부 AUC, 학습 횟수 민감도, 학습 곡선은 이후에 추가한 탐색적 분석으로 구분하여 보고한다.
신경망 모델은 모든 코호트에서 seed 42로 학습하고, 가장 큰 두 코호트에서 seed 13과 2026을 추가하여 학습 초기화에 따른 변동을 확인한다.
seed 변동은 두 코호트에서만 확인하였으므로, 작은 코호트의 결과는 이 한계 안에서 해석한다.
학습 횟수의 상한이 신경망 모델의 성능을 제한하는지 확인하기 위해, 같은 두 코호트에서 상한을 24 epoch로 늘려 다시 학습한다.
데이터 규모의 영향은 같은 두 코호트에서 학습 경기 수를 300, 600, 1,200경기와 전체로 바꾼 학습 곡선으로 확인하며, epoch 수를 고정한 경우와 갱신 횟수를 맞춘 경우를 모두 구한다.


\bibliographystyle{IEEEtran}
\bibliography{review}
\end{document}
